\expandafter\ifx\csname CAPETINCLUDEMODE\endcsname\relax
\documentclass[11pt,a4paper]{article}
% preamble.tex — packages et commandes communes a toutes les fiches CAPET SII IE
\input{glyphtounicode}
\pdfgentounicode=1
\usepackage{cmap}
\usepackage[T1]{fontenc}
\usepackage[utf8]{inputenc}
\usepackage[french]{babel}
\usepackage[a4paper, top=1.6cm, bottom=1.6cm, left=1.4cm, right=1.4cm]{geometry}
\setlength{\headheight}{14.5pt}
\usepackage{amsmath, amssymb, amsfonts}
\usepackage{siunitx}
\sisetup{locale=FR, output-decimal-marker={,}, inter-unit-product=\ensuremath{{}\cdot{}}}
\usepackage{array, booktabs, multirow, tabularx}
\usepackage{xcolor}
\usepackage[most, breakable]{tcolorbox}
\usepackage{multicol}
\usepackage{enumitem}
\usepackage{mathtools}
\usepackage{lmodern}
\usepackage{microtype}
\usepackage{fancyhdr}
\usepackage{tikz}
\usetikzlibrary{positioning,arrows.meta}
\usepackage{lastpage}

% ── Couleurs par domaine ──────────────────────────────────────────────
\definecolor{colElec}{RGB}{26, 82, 148}       % Électrotechnique — bleu foncé
\definecolor{colElecLight}{RGB}{210, 225, 245} % fond bleu clair
\definecolor{colAuto}{RGB}{34, 120, 66}        % Automatique — vert
\definecolor{colMeca}{RGB}{185, 90, 20}        % Mécanique — orange
\definecolor{colTherm}{RGB}{160, 40, 40}       % Thermique — rouge
\definecolor{colQT}{RGB}{232, 244, 233}        % fond question type (vert très léger)
\definecolor{colJury}{RGB}{255, 243, 205}      % fond jury (jaune ambre)
\definecolor{colJuryFrame}{RGB}{200, 130, 0}   % bordure jury
\definecolor{colGrayLight}{RGB}{245, 245, 247} % fond gris neutre
\definecolor{colSeparator}{RGB}{180, 180, 180} % gris séparateur

% ── Style tcolorbox de section principale ────────────────────────────
\tcbset{
  ficheSection/.style={
    breakable,
    enhanced,
    colback=#1!7!white,
    colframe=#1,
    colbacktitle=#1,
    coltitle=white,
    fonttitle=\bfseries\normalsize,
    left=4pt, right=4pt, top=3pt, bottom=4pt,
    before skip=6pt, after skip=6pt,
  }
}

% ── Style encadré question type ──────────────────────────────────────
\tcbset{
  questiontype/.style={
    breakable,
    enhanced,
    colback=colQT,
    colframe=colElec!55!black,
    colbacktitle=colElec!55!black,
    coltitle=white,
    fonttitle=\bfseries\small,
    left=4pt, right=4pt, top=3pt, bottom=3pt,
    before skip=5pt, after skip=5pt,
  }
}

% ── Style encadré jury ───────────────────────────────────────────────
\tcbset{
  jury/.style={
    breakable,
    enhanced,
    colback=colJury,
    colframe=colJuryFrame,
    fonttitle=\bfseries\normalsize,
    left=4pt, right=4pt, top=3pt, bottom=3pt,
  }
}

% ── Mise en page ─────────────────────────────────────────────────────
\pagestyle{fancy}
\fancyhf{}
\renewcommand{\headrulewidth}{0.4pt}
\renewcommand{\footrulewidth}{0.3pt}

\setlist[itemize]{leftmargin=1.2em, itemsep=1pt, topsep=2pt, parsep=0pt}
\setlist[enumerate]{leftmargin=1.2em, itemsep=3pt, topsep=2pt, parsep=0pt}

\setlength{\parskip}{2pt}
\setlength{\columnsep}{10pt}

\def\CAPETBEGINFILE{\begin{document}}
\def\CAPETENDFILE{\end{document}}
\else
\def\CAPETBEGINFILE{}
\def\CAPETENDFILE{}
\fi

\fancyhead[L]{\textbf{\textcolor{colElec}{T25}} - Transmission numerique - Modulations et protocoles}
\fancyhead[R]{\small\textit{CAPET SII IE - Session 2026}}
\fancyfoot[C]{\small\thepage\,/\,\pageref{LastPage}}

\CAPETBEGINFILE

% -----------------------------------------------------------------------
%  BLOC EN-TeTE
% -----------------------------------------------------------------------

\begin{tcolorbox}[
  enhanced, breakable,
  colback=colElec,
  colframe=colElec,
  left=6pt, right=6pt, top=5pt, bottom=5pt,
  before skip=0pt, after skip=4pt
]
  {\color{white}\bfseries\large T25 \quad Transmission numerique - Modulations et protocoles}\\[2pt]
  {\color{colElec!30!white}\small
    \textbf{Domaine :} electronique - Communications numeriques \quad
    \textbf{Presence :} \textbf{6/13} sujets (2013-2026) - theme \emph{frequent}
  }\\[1pt]
  {\color{colElec!20!white}\footnotesize
    Sous-themes : debit binaire - modulation FSK/FM - LoRa/LoRaWAN - bus CAN (trame) - MLI/PWM transmission - Shannon capacite - UART vs PWM
  }
\end{tcolorbox}

% -----------------------------------------------------------------------
%  1. RAPPELS DE COURS
% -----------------------------------------------------------------------

\begin{tcolorbox}[
  ficheSection=colElec,
  title={1. Rappels de cours},
]

% -- 1.1 Debit binaire -------------------------------------------------
\textbf{1.1 Notions fondamentales}

\[
  D_b = \frac{1}{T_b} \quad [\mathrm{bit/s}]
  \qquad
  D_b = \log_2(M) \cdot D_s \quad (M \text{ etats})
  \qquad
  C = B \cdot \log_2(1 + SNR) \quad \text{(Shannon)}
\]

\medskip\hrule\vspace{4pt}

% -- 1.2 Modulation FM et FSK ------------------------------------------
\textbf{1.2 Modulation FM et FSK (2015 E1)}

\[
  f_{inst}(t) = f_0 + k_f\,v_m(t)
  \qquad
  m = \frac{\Delta f}{f_m} = \frac{k_f\,V_m}{f_m} \quad \text{(indice de modulation)}
\]

\textbf{FSK} ($v_m = \pm 1$) :
\begin{itemize}
  \item Bit - 0 - : frequence $f_0 = f_n - \Delta f$
  \item Bit - 1 - : frequence $f_1 = f_n + \Delta f$
  \item Excursion : $\Delta f = f_1 - f_n = f_n - f_0$
\end{itemize}
Exemple (2015) : $f_n = 433\,\mathrm{MHz}$, $\Delta f = 300\,\mathrm{kHz}$.

\medskip\hrule\vspace{4pt}

% -- 1.3 LoRa / LoRaWAN ------------------------------------------------
\textbf{1.3 LoRa / LoRaWAN (2022 E1/E2)}

Modulation CSS (Chirp Spread Spectrum), bandes ISM 868 MHz (Europe), 915 MHz (USA).
\begin{center}
\renewcommand{\arraystretch}{1.2}
\begin{tabular}{ccccc}
\toprule
\textbf{DR} & \textbf{Mod.} & \textbf{SF} & \textbf{BP (kHz)} & \textbf{Debit physique (bit/s)} \\
\midrule
0 & LoRa & 12 & 125 & 250 \\
3 & LoRa & 9  & 125 & 1\,760 \\
5 & LoRa & 7  & 125 & 5\,470 \\
7 & FSK  & -  & -   & 50\,000 \\
\bottomrule
\end{tabular}
\end{center}
Debit utile avec taux de codage CR = 4/5 : $D_{utile} = D_{phys} \times 4/5$.\\
Portee : $\sim 5\,\mathrm{km}$ urbain, $\sim 15\,\mathrm{km}$ rural (SF12).

\medskip\hrule\vspace{4pt}

% -- 1.4 Bus CAN -------------------------------------------------------
\textbf{1.4 Bus CAN - transmission serie de terrain}

Bus serie differentiel ($U_D = U_{CAN\_H} - U_{CAN\_L}$), multi-maetres, 1 Mbit/s max.

Structure de trame (donnees = 8 octets) : 111 bits au total.
\[
  \underbrace{SOF}_{1} + \underbrace{ID}_{11} + \underbrace{RTR}_{1} + \underbrace{IDE+r0}_{2} + \underbrace{DLC}_{4} + \underbrace{Data}_{0-64} + \underbrace{CRC}_{15+1} + \underbrace{ACK}_{2} + \underbrace{EOF}_{7}
\]

\textbf{Debit depuis chronogramme :} mesurer $T_b$ sur le chrono - $D_b = 1/T_b$ ; duree trame $= 111 \times T_b$.

\medskip\hrule\vspace{4pt}

% -- 1.5 MLI/PWM en transmission ---------------------------------------
\textbf{1.5 MLI/PWM - contexte transmission de consigne}

\begin{center}
\renewcommand{\arraystretch}{1.2}
\begin{tabular}{lll}
\toprule
\textbf{Critere} & \textbf{UART} & \textbf{PWM} \\
\midrule
Synchronisation & Asynchrone (baud rate) & Synchrone (frequence fixe) \\
Complexite & elevee (protocole complet) & Simple (mesure $t_{on}$) \\
Immunite bruit & Faible & elevee (signal TOR) \\
Debit & eleve & Faible (1 valeur/periode) \\
\bottomrule
\end{tabular}
\end{center}

\end{tcolorbox}

% -----------------------------------------------------------------------
%  2. FORMULAIRE SYNTHeTIQUE
% -----------------------------------------------------------------------

\begin{tcolorbox}[
  ficheSection=colElec,
  title={2. Formulaire synthetique},
]
\small
\begin{multicols}{2}
\renewcommand{\arraystretch}{1.45}
\begin{tabular}{@{}p{5cm}ll@{}}
\toprule
\textbf{Formule} & \textbf{Usage} & \textbf{U.} \\
\midrule
\multicolumn{3}{@{}l}{\footnotesize\textit{Debit}} \\
$D_b = 1/T_b$ & Bit/s & bit/s \\
$D_b = \log_2(M) \cdot D_s$ & $M$ etats & bit/s \\
$C = B\log_2(1+SNR)$ & Shannon & bit/s \\
\midrule
\multicolumn{3}{@{}l}{\footnotesize\textit{FM / FSK}} \\
$f_{inst} = f_0 + k_f v_m$ & Inst. & Hz \\
$m = \Delta f / f_m$ & Indice & - \\
$f_1 = f_n + \Delta f$ ; $f_0 = f_n - \Delta f$ & FSK & Hz \\
\bottomrule
\end{tabular}

\columnbreak

\renewcommand{\arraystretch}{1.45}
\begin{tabular}{@{}p{5cm}ll@{}}
\toprule
\textbf{Formule} & \textbf{Usage} & \textbf{U.} \\
\midrule
\multicolumn{3}{@{}l}{\footnotesize\textit{LoRa}} \\
$D_{utile} = D_{phys} \times 4/5$ & CR=4/5 & bit/s \\
SF$\uparrow$ - debit$\downarrow$, portee$\uparrow$ & Compromis & - \\
\midrule
\multicolumn{3}{@{}l}{\footnotesize\textit{Bus CAN}} \\
Debit max : 1\,Mbit/s & CAN 2.0 & bit/s \\
Trame max : 111 bits & - & bits \\
$t_{trame} = 111 \times T_b$ & Duree & s \\
\midrule
\multicolumn{3}{@{}l}{\footnotesize\textit{Donnees types}} \\
LoRa SF7/125 kHz : 5\,470 bit/s & DR5 & - \\
FSK LoRa : 50\,000 bit/s & DR7 & - \\
\bottomrule
\end{tabular}
\end{multicols}
\normalsize
\end{tcolorbox}

% -----------------------------------------------------------------------
%  3. QUESTIONS TYPES
% -----------------------------------------------------------------------

\begin{tcolorbox}[
  ficheSection=colElec,
  title={3. Questions types - extraites des sujets reels},
]

% -- QT1 -------------------------------------------------------------

% Mini-procedure de resolution type copie concours
\begin{tcolorbox}[questiontype,title={Procedure type copie concours (5 etapes)}]
\begin{enumerate}[leftmargin=1.5em]
  \item \textbf{Qualifier le modele} : hypotheses, domaine de validite, grandeurs d'entree/sortie.
  \item \textbf{Ecrire les lois} : equations de base avec symboles et unites SI.
  \item \textbf{Isoler l'inconnue} : expression litterale, puis application numerique tracee.
  \item \textbf{Controler l'ordre de grandeur} : coherence physique et unite finale.
  \item \textbf{Conclure concours} : phrase courte orientee decision technique.
\end{enumerate}
\end{tcolorbox}

% Visuel de methode : chemin de resolution CAPET
\begin{center}
\begin{tikzpicture}[font=\small, node distance=7mm and 6mm, >=Stealth]
  \node[draw, rounded corners, fill=colElec!8, inner sep=4pt] (n1) {Notions};
  \node[draw, rounded corners, fill=colElec!8, right=of n1, inner sep=4pt] (n2) {Methode};
  \node[draw, rounded corners, fill=colElec!8, right=of n2, inner sep=4pt] (n3) {Application annale};
  \node[draw, rounded corners, fill=colElec!8, right=of n3, inner sep=4pt] (n4) {Verification jury};
  \draw[->, thick, colElec] (n1) -- (n2);
  \draw[->, thick, colElec] (n2) -- (n3);
  \draw[->, thick, colElec] (n3) -- (n4);
\end{tikzpicture}
\end{center}
\begin{tcolorbox}[questiontype,
  title={QT1 - Modulation FSK - deux frequences et indice \hfill\normalfont\footnotesize E1 2015 (Shelter deployable)}
]
\textbf{enonce.}
Systeme de transmission FSK : porteuse $f_n = 433\,\mathrm{MHz}$, excursion $\Delta f = 300\,\mathrm{kHz}$, frequence modulante $f_m = 1\,\mathrm{kHz}$. (a) Donner les deux frequences FSK. (b) Calculer l'indice de modulation $m$. (c) Exprimer la phase instantanee $\Phi(t)$.

\tcbline

\textbf{Corrige.}

(a) Frequences FSK :
\[
  f_1 = f_n + \Delta f = 433 + 0{,}3 = 433{,}3\,\mathrm{MHz} \quad (\text{bit e1e})
\]
\[
  f_0 = f_n - \Delta f = 433 - 0{,}3 = 432{,}7\,\mathrm{MHz} \quad (\text{bit e0e})
\]

(b) Indice de modulation :
\[
  m = \frac{\Delta f}{f_m} = \frac{300 \times 10^3}{1000} = \mathbf{300}
\]

(c) Phase instantanee pour $v_m(t) = V_m \sin(2\pi f_m t)$ :
\[
  \Phi(t) = 2\pi f_n t + m \sin(2\pi f_m t) = 2\pi f_n t + 300\sin(2\pi \times 10^3 t)
\]
\end{tcolorbox}

\vspace{4pt}

% -- QT2 -------------------------------------------------------------

\begin{tcolorbox}[questiontype,
  title={QT2 - LoRa - debit utile et choix du SF \hfill\normalfont\footnotesize E1/E2 2022 (Q29)}
]
\textbf{enonce.}
Reseau LoRaWAN avec SF = 7, BP = 125 kHz, CR = 4/5. (a) Lire le debit physique dans la table. (b) Calculer le debit binaire utile. (c) Comparer SF7 et SF12 pour une application IoT necessitant portee maximale.

\tcbline

\textbf{Corrige.}

(a) Table : SF7 / 125 kHz - $D_{phys} = 5\,470\,\mathrm{bit/s}$.

(b) Debit utile :
\[
  D_{utile} = 5\,470 \times \frac{4}{5} = \mathbf{4\,376\,\mathrm{bit/s}}
\]

(c) SF12 / 125 kHz - $D_{phys} = 250\,\mathrm{bit/s}$ seulement, mais portee $\sim 15\,\mathrm{km}$ en zone rurale.\\
Compromis : SF12 pour longue portee (objets isoles) ; SF7 pour debit maximal (zone urbaine dense).
\end{tcolorbox}

\vspace{4pt}

% -- QT3 -------------------------------------------------------------

\begin{tcolorbox}[questiontype,
  title={QT3 - Bus CAN - debit depuis chronogramme \hfill\normalfont\footnotesize E1 2026 (D5)}
]
\textbf{enonce.}
Sur un chronogramme du bus CAN, la duree d'un bit est $T_b = 1\,\mu\mathrm{s}$. Calculer le debit binaire. En deduire la duree d'une trame complete (8 octets de donnees). Verifier le cahier des charges (trame $< 200\,\mu\mathrm{s}$).

\tcbline

\textbf{Corrige.}

Debit binaire :
\[
  D_b = \frac{1}{T_b} = \frac{1}{10^{-6}} = \mathbf{1\,\mathrm{Mbit/s}}
\]

Duree de la trame (111 bits avec 8 octets de donnees) :
\[
  t_{trame} = 111 \times T_b = 111 \times 1\,\mu\mathrm{s} = \mathbf{111\,\mu\mathrm{s}}
\]

Conclusion : $111\,\mu\mathrm{s} < 200\,\mu\mathrm{s}$ - cahier des charges respecte .
\end{tcolorbox}

\end{tcolorbox}

% -----------------------------------------------------------------------
%  4. ATTENTES DU JURY
% -----------------------------------------------------------------------

\begin{tcolorbox}[
  jury,
  title={4. Attentes du jury - Erreurs recurrentes et points de vigilance},
]

\begin{enumerate}[label=\textbf{\arabic*.}, leftmargin=*, itemsep=5pt]

  \item \textbf{Confondre bus CAN et CAN (convertisseur).} Meme sigle, deux significations distinctes. Dans les sujets 2026, les deux coexistent : se fier au contexte (bus de terrain - peripherique microcontroleur).

  \item \textbf{FSK est un cas particulier de FM.} Modulant binaire $v_m = \pm 1$ - deux frequences fixes. Relier la formule generale $f_{inst} = f_0 + k_f v_m$ au cas FSK.

  \item \textbf{Taux de codage LoRa.} CR = 4/5 signifie 4 bits utiles + 1 bit redondant = 5 bits transmis. Debit utile = $D_{phys} \times 4/5$. Ne pas confondre CR et efficacite spectrale.

  \item \textbf{Chronogramme CAN.} Mesurer $T_b$ en comptant les transitions, pas les niveaux. Une erreur de lecture sur $T_b$ se propage au debit et - la duree de trame.

  \item \textbf{Shannon (capacite) vs Shannon (echantillonnage).} $C = B\log_2(1+SNR)$ est la capacite maximale d'un canal bruite. $f_e \geq 2f_{max}$ est le theoreme d'echantillonnage - deux resultats distincts du meme auteur.

  \item \textbf{MLI onduleur vs PWM servomoteur.} La MLI de commande d'onduleur (T09) compare une modulante sinusoedale - une porteuse triangulaire. La PWM de consigne servo (50 Hz, $t_{on}$ proportionnel - l'angle) est un protocole de transmission - ne pas confondre.

\end{enumerate}

\medskip
\begin{tcolorbox}[
  colback=colJuryFrame!12!colJury,
  colframe=colJuryFrame,
  left=4pt, right=4pt, top=2pt, bottom=2pt,
  fontupper=\small\itshape
]
Rapport jury 2022 : "La partie LoRa a ete bien traitee par les candidats qui avaient lu la documentation. Le calcul du debit utile doit imperativement tenir compte du taux de codage CR, souvent oublie."\\
Rapport jury 2026 : "La structure d'une trame CAN doit etre connue. La lecture de chronogramme est une competence instrumentale de base attendue du futur enseignant."
\end{tcolorbox}

\end{tcolorbox}

% -----------------------------------------------------------------------
%  PIED - PReSENCE DANS LES SUJETS
% -----------------------------------------------------------------------

\begin{tcolorbox}[
  enhanced, breakable,
  colback=colGrayLight,
  colframe=colSeparator,
  left=4pt, right=4pt, top=3pt, bottom=3pt,
  before skip=4pt, after skip=0pt,
  fonttitle=\small\bfseries,
  title={Presence dans les sujets analyses (2013-2026)},
]
\begin{center}
\footnotesize\renewcommand{\arraystretch}{1.15}
\begin{tabular}{lccccccccccccc}
\toprule
\textbf{Annee} & 13 & 14 & 15 & 16 & 17 & 18 & 19 & 20 & 21 & 22 & 23 & 24 & 25 \\
\midrule
E1 (commune)   & - & - & $\checkmark$ & - & - & - & $\checkmark$ & - & - & $\checkmark$ & - & - & - \\
E2 discipl. IE & - & - & - & - & - & - & $\checkmark$ & - & - & - & - & - & $\checkmark$ \\
E2 appl. comm. & - & - & - & - & - & - & - & - & - & - & - & - & $\checkmark$ \\
\bottomrule
\end{tabular}
\end{center}
\smallskip
\footnotesize\textbf{Sous-themes frequents :}
modulation FSK (2015) - LoRa SF/BP/CR debit (2022) - bus CAN trame + chronogramme (2019, 2026) - PWM vs UART (2025) - debit binaire depuis chronogramme

\medskip
\textbf{Ordres de grandeur (reperes rapides).}
\begin{itemize}[leftmargin=1.3em]
  \item Reseau BT FR : tension phase-neutre typique \SI{230}{\volt}, tension phase-phase typique \SI{400}{\volt}, frequence \SI{50}{\hertz}.
  \item Courants d'etude frequents en sujet : de quelques \si{\ampere} a quelques dizaines de \si{\ampere}.
  \item En conversion d'energie, un rendement attendu reste generalement entre 0{,}8 et 0{,}98 selon la technologie.
\end{itemize}

\medskip
\textbf{Synthese express (a refaire sans support).}
\begin{enumerate}[leftmargin=1.4em]
  \item Poser les hypotheses du modele et les conventions de signe.
  \item Ecrire 2 a 4 relations cles avec unites.
  \item Realiser un calcul numerique complet et verifier l'ordre de grandeur.
  \item Justifier le resultat par un argument technique de niveau CAPET.
\end{enumerate}

\end{tcolorbox}

\CAPETENDFILE









